%%%PAGE 102 rot=0 legibility=good summary=实验：用单摆测重力加速度（原理、数据处理、误差、注意事项）；实验：插针法测玻璃折射率（图像法、单位圆法、误差、注意事项）
%%%BLOCK id=102-01 ch=08 sec=用单摆测重力加速度 kind=method exp=1 alt=none cont=none y=0.03-0.58
\textbf{探究单摆的运动\quad 用单摆测定重力加速度}

\textbf{实验原理与操作}
\[ T=2\pi\sqrt{\frac{L}{g}}\qquad\qquad g=\frac{4\pi^{2}L}{T^{2}} \]

\textbf{数据处理}
\begin{itemize}
\item 公式法
\item 图像法.
\end{itemize}

\textbf{误差分析}
\begin{itemize}
\item 偶然误差.
 \begin{itemize}
 \item 测量时间及摆长时产生的误差
 \item 减小方法：多次测量求平均值/从单摆经过平衡位置时开始计时
 \end{itemize}
\item 系统误差
 \begin{itemize}
 \item 单摆自身
 \item 摆球选体积小，密度大的
 \item 最大摆角要小于$5^\circ$.
 \end{itemize}
\end{itemize}

\textbf{注意事项}
\begin{itemize}
\item 用$1\unit{m}$左右细线
\item 悬线顶端不能晃动用夹子夹住，保证顶点固定
\item 在同一竖直面内摆动，摆角小于$5^\circ$. 从摆球摆到平衡位置开始计时，并数准全振动的次数.
\end{itemize}
%%%END
%%%BLOCK id=102-02 ch=15 sec=测玻璃折射率·插针法 kind=method exp=1 alt=none cont=none y=0.56-1.00
\textbf{测定玻璃折射率\quad 插针法}

\textbf{实验原理与操作}\quad 铺白纸画线、插针、测量
\[ n=\frac{\sin\theta_1}{\sin\theta_2} \]

\textbf{数据处理}
\begin{itemize}
\item 计算法\quad 计算每次折射率，求出平均值$\bar{n}$
\item 图像法
\end{itemize}

%FIG p102_a 0.37 0.722 0.54 0.795
\notefig[0.28]{figures/p102_a.png}

$\sin\theta_1$—$\sin\theta_2$图像\qquad $n=k=\dfrac{\sin\theta_1}{\sin\theta_2}$

\textbf{误差分析}\quad 单位圆法

%FIG p102_b 0.325 0.785 0.495 0.92
\notefig[0.3]{figures/p102_b.png}

\[ n=\frac{EH}{E'H'} \]

\begin{itemize}
\item[\unclear{①}] 根据插针位置画出入射光线，出射光线时有偏差. % 圈号①位于页边被截断，据②③推测
\item[②] 测量时入射角与折射角的测量误差.
\item[③] 作图时，玻璃砖的宽度绘制有误差，画的两边界不平行出现误差.
\end{itemize}

\textbf{注意事项}
\begin{itemize}
\item 玻璃砖要用厚度大的
\item 入射角在$30^\circ\sim60^\circ$之间
\item 大头针要竖直插在白纸上，且距离应尽量大一些
\item 玻璃砖的折射面要画准，不能用玻璃砖界面代替直尺画界线
\end{itemize}
%%%END
%%%PAGEEND 102
